

\section{Formulación del Problema}
\label{sc:FP}

Habiendo revisado varios trabajos en la literatura, se observa que continúa la investigación de diversos métodos para mejorar los algoritmos de optimización con diferentes enfoques. A pesar de los avances, siempre se enfrenta el problema del estancamiento en las metaheurísticas.


\subsection{El Problema del Estancamiento en Metaheurísticas}

El estancamiento en las metaheurísticas es un fenómeno que ocurre cuando un algoritmo de optimización, que busca encontrar la solución óptima o una solución cercana al óptimo en un espacio de búsqueda, se queda atrapado en una región del espacio de soluciones y es incapaz de encontrar mejores soluciones a pesar de seguir iterando, quedando estancado en lo que se llama un óptimo local. Este problema se presenta comúnmente en los algoritmos bio-inspirados que resuelven problemas de optimización combinatorial, donde el espacio de búsqueda es vasto y complejo \cite{10.5555/1718024, Michalewicz2004}.

\subsubsection{Características del Estancamiento}

El estancamiento se caracteriza por:
\begin{itemize}
    \item \textbf{Convergencia Prematura:} El algoritmo converge rápidamente a una solución subóptima y no explora adecuadamente otras regiones del espacio de búsqueda.
    \item \textbf{Falta de Diversidad:} La población de soluciones en algoritmos basados en población (como los algoritmos genéticos) se vuelve homogénea, reduciendo la capacidad de encontrar nuevas soluciones.
    \item \textbf{Exploración Insuficiente:} El algoritmo explora insuficientemente el espacio de búsqueda, centrando sus esfuerzos en áreas ya conocidas y dejando de lado regiones potencialmente prometedoras.
\end{itemize}

\subsubsection{Causas del Estancamiento}

Según \cite{Dokeroglu2019_jx} y de acuerdo a lo revisado en la literatura, las principales causas del estancamiento incluyen:
\begin{itemize}
    \item \textbf{Función de Evaluación Engañosa:} La función objetivo puede tener múltiples óptimos locales que engañan al algoritmo para que se quede atrapado en una de estas soluciones.
    \item \textbf{Parámetros del Algoritmo:} Parámetros mal ajustados, como tasas de mutación y cruce en algoritmos genéticos u otros, pueden llevar a una exploración insuficiente del espacio de búsqueda.
    \item \textbf{Estrategia de Búsqueda:} Estrategias de búsqueda que favorecen la explotación sobre la exploración pueden aumentar la probabilidad de estancamiento.
\end{itemize}

\subsubsection{Consecuencias del Estancamiento}

El estancamiento puede llevar a:
\begin{itemize}
    \item \textbf{Soluciones Subóptimas:} El algoritmo puede terminar con soluciones que están lejos del óptimo global.
    \item \textbf{Tiempo de Cómputo Elevado:} Más iteraciones sin mejora significativa resultan en un uso ineficiente de recursos computacionales.
\end{itemize}

\subsubsection{Métodos para Superar el Estancamiento}

Para superar el estancamiento, en el estado del arte se pudo encontrar varias maneras, y se pueden aplicar diversas estrategias como:
\begin{itemize}
    \item \textbf{Métodos Híbridos:} Combinar diferentes algoritmos de optimización para aprovechar sus fortalezas.
    \item \textbf{Técnicas de Diversificación:} Introducir mecanismos que aumenten la diversidad de la población de soluciones.
    \item \textbf{Reinicialización:} Reiniciar el algoritmo desde una nueva posición en el espacio de búsqueda.
    \item \textbf{Ajuste Dinámico de Parámetros:} Modificar los parámetros del algoritmo durante la ejecución para mejorar la exploración y explotación.
\end{itemize}

Varios de estos métodos utilizan distintos enfoques que manipulan los parámetros de las metaheurísticas. En este trabajo se propone otra manera de resolver el problema del estancamiento utilizando un modelo estadístico bayesiano que pueda ser explicado por SHAP para salir del estancamiento ingeniosamente según las métricas que nos entregue el algoritmo en cuestión.










\section{Solución Propuesta}
\label{sc:SP}

Considerando la falta de trabajos realizados en el contexto de metaheurísticas que sean explicables, en este estudio se propone un método híbrido hacia algoritmos de optimización que les permita utilizar las ventajas de entrenar a un modelo bayesiano durante el proceso de optimización (online), como un modelo que predice el fitness, en conjunto a un método de detección del estancamiento de las soluciones, y que mediante herramientas de explicabilidad como SHAP, se pueda determinar cuál esta siendo el factor causante del estancamiento a través del modelo bayesiano, lo que permitirá actualizar las soluciones de una manera ingeniosa y enfocada al problema.


\begin{figure}[ht!]
    \centering
    \includegraphics[width=0.6\textwidth]{imagenes/diagramaFlujoPropusta.png}
    \caption{Diagrama de flujo de la propuesta de solución.}
    \label{fig:dflujo_propuesta}
\end{figure}

La propuesta se ve en la Figura \ref{fig:dflujo_propuesta}. Se presenta la estructura general de una metaheurística, la que abarca varias implementaciones de algoritmos bio-inspirados. Se comienza con la inicialización de las variables, la población inicial de soluciones y el fitness de estas.

Para esta propuesta, cada uno de los individuos de la población tendrá un Modelo Bayesiano Regresivo (MBR), al cual se le entregará el vector solución y el fitness calculado por cada iteración, a modo de que se entrene de manera "online", a medida que datos nuevos van surgiendo. Esto es lo que representa el color verde del diagrama.

Por otra parte, el algoritmo poseerá un método de detección de estancamiento, plasmada en color naranjo dentro de la figura. En especial, para este estudio se utilizará el ``Monitoreo del cálculo de la Mejora Relativa'' (MMR). Este funcionará de manera individual, es decir, irá tomando el historial de los mejores fitness que se han actualizado de este agente para determinar después de la 1/4 parte del criterio de parada (ventana de detección) si el individuo se encuentra estancado o no. Para calcular la MMR se utiliza la siguiente fórmula:

La mejora relativa se define como:

\[
\text{Mejora Relativa} = \frac{f_1 - f_n}{f_1}
\label{ec:mejora_relativa_1}
\]

La condición para determinar si ha habido una mejora significativa es:

\[
\frac{f_1 - f_n}{f_1} > \epsilon
\label{ec:mejora_relativa_2}
\]

donde \( f_1 \) es el valor de fitness al inicio de la ventana de detección, \( f_n \) es el valor de fitness en la iteración más reciente dentro de esa ventana, y \(\epsilon\) es el umbral de mejora, que preliminarmente se le asignará un valor de $0.01$.

Por último, los procesos de color azul son la parte de la sección de explicabilidad (XAI) del comportamiento de las variables de decisión con respecto al fitness. Este proceso será utilizado solo en caso de que el individuo haya sido detectado en estancamiento. Para explicar las variables que están influyendo se utiliza el MBR propio del agente, que posee el historial de datos que ha utilizado para ir actualizando el fitness. La técnica empleada en esta sección es SHAP, que permitirá ver cuales son las variables que menos están aportando, con el fin de actualizar al individuo con las variables que realmente están aportando más cambio al fitness (maximización o minimización).

El método para salir del estancamiento, se realiza con la selección del número de variables que serán tomadas en cuenta de SHAP, que dependerán de la probabilidad de crossover o mutación que exista, y principalmente del porcentaje de importancia que tenga unas variables respecto a otras.

Una vez, se ejecuta el método para salir del estancamiento, se reentrena el MBR, se comparan las soluciones, se actualiza la mejor solución global de la población, y se verifica el criterio de parada para continuar con el ciclo o finalizarlo. El pseudocódigo de este diagrama puede ser visto en detalle \ref{alg:metaheuristics_bayesianXAI}.



\begin{algorithm}[!ht]
    \scriptsize
    \caption{Estructura común de las metaheurísticas con la propuesta en peudocódigo}\label{alg:metaheuristics_bayesianXAI}
    \KwIn{datos de entrada del problema, $popSize$, $maxIter$, $detectionWindow$, $improvement_threshold$}
    \KwOut{una solución óptima obtenida durante el proceso de optimización.}

    cargarParámetrosDelAlgoritmo()\;
    cargarDatosDelProblema()\;
    \tcp{el valor $n$ define el número de dimensiones}
    Define la función objetivo $f(x)$, $x=\langle x^{1},\ldots,x^{n} \rangle$\;
    \tcp{produce la primera generación de $popSize$ agentes}
    \ForEach{agente $a_{i}$ \textbf{desde} 1 \textbf{hasta} $popSize$}{
        \ForEach{dimensión $j$ \textbf{desde} 1 \textbf{hasta} $n$}{
            \While{posición $x_{i}^{j}$  no sea factible}{
                posición $x_{i}^{j} \leftarrow$ ValorAleatorioEntre\{0, 1\}\;
                Inicialización de variables adicionales de la metaheurística\;
            }
        }
        Maximizar o minimizar la función objetivo y comparar para obtener la mejor solución inicial\;
        $MBR\_x_{i}$ $\leftarrow$ Entrenar Modelo Bayesiano Regresivo del agente $i$\;
    }
    \tcp{producir $maxIter$-generaciones de $popSize$ agentes}
    \While{$t < maxIter$}{
        \tcp{fase del proceso es invocado}
        Asigna valores de diversificación o intensificación (dependiendo de $t$)\;
        \ForEach{agente $a_{i}$ \textbf{desde} 1 \textbf{hasta} $popSize$}{
            \ForEach{dimensión $j$ \textbf{desde} 1 \textbf{hasta} $n$}{
                \tcp{genera nuevas soluciones}
                \While{posición $x_{i}^{j}$ no es factible}{
                    posición $x_{i}^{j} \leftarrow$ es actualizada por las ecuaciones de la metaheuristica\;
                }
            }
            Maximizar o minimizar la función objetivo (calculo del fitness)\;

            actualizar $MBR\_x_{i}$ entrenándolo en línea con nueva data del agente $i$\;

            \If{($t > detectionWindow$)}{
                MMR $\leftarrow$ calcular método de monitoreo de la mejora relativa de la Ec. \ref{ec:mejora_relativa_1} y \ref{ec:mejora_relativa_2}\;
                \If{$MMR$ < $improvement_threshold$}{
                    ejecutar método XAI a través de SHAP para explicar $MBR$ de la solución $i$\;
                    $arrayResponsable$ $\leftarrow$ seleccionar alguna(s) variable(s) responsable(s) estancamiento\;
                    \If{randomValorEntre(0,1) > $probMutacion$}{
                        Ejecutar mutación en $arrayResponsable$\;
                    }
                    \Else{
                        \If{randomValorEntre(0,1) > $probCrossover$}{
                            Ejecutar cruce (crossover) en $arrayResponsable$\;
                        }
                    }
                    Maximizar o minimizar la función objetivo (calculo del fitness)\;

                    actualizar $MBR\_x_{i}$ entrenándolo en línea con nueva data del agente $i$\;
                }
            }
            actualizar la mejor solución \;
        }
    }
    Retornar resultados y visualización post-proceso\;
\end{algorithm}

La propuesta de mejora a las metaheurísticas en el ámbito académico, abre nuevas áreas de investigación y aplicación, fomentando la innovación y el desarrollo de nuevas técnicas y algoritmos. Esto puede llevar a la formación de nuevas teorías y modelos, enriqueciendo el conocimiento en el campo de la optimización. Además, puede facilitar la resolución de problemas complejos de la vida real de manera ingeniosa en diferentes disciplinas.

Además puede tener un impacto significativo en varios ámbitos si es que se aplica en la problemas reales. En el ámbito comercial, empresas que dependen de la optimización de recursos y procesos, como las de logística y manufactura, pueden beneficiarse enormemente de soluciones más eficientes. Esto podría resultar en una reducción de costos operativos, mejor utilización de recursos y una mayor competitividad en el mercado. En el ámbito social, mejoras en la eficiencia de procesos industriales y de transporte pueden contribuir a una reducción de la huella de carbono, promoviendo prácticas más sostenibles y amigables con el medio ambiente.







\section{Pregunta de Investigación}
\label{sc:PI}

Basándonos en la revisión de la literatura, se plantean las siguientes preguntas de investigación:

%Desarrolla al menos una pregunta de investigación, claramente relevante para la validación del tratamiento. (estas podrán ser respondidas luego de haber ejecutado los casos experimentales definidos)%

\begin{itemize}
    \item ¿Cuál es la diferencia de la calidad del algoritmo propuesto en comparación a las versiones originales de los algoritmos de optimización implementados en la resolución de funciones benchmark y un problema real para un mismo contexto?. %esta pregunta nos permite averiguar qué tan bueno es el algoritmo dada una instancia en particular

    \item ¿Cuál es el impacto de la escalabilidad en la calidad de las soluciones del algoritmo con la mejora propuesta respecto al algoritmo en su versión original?, es decir, en instancias más grandes del problema. %esta pregunta nos permite saber si la instancia funciona mejor cuando la cantidad de puntos de entrega del problema es pequeño o grande

\end{itemize}










\section{Hipótesis}
\label{sc:HIPO}

%Define hipótesis nula y alternativa consistentes con la pregunta de investigación y las mediciones a realizar.%
Para esta investigación se planean hipótesis según el diseño experimental que se implementa. Debido a que la comparación es pareada entre el algoritmo con la propuesta y en su versión original, se utilizarán dos hipótesis para cada par de algoritmos que se utilice en este trabajo, las cuales son:
\begin{itemize}
    \item Hipótesis nula (H0): ``No hay diferencia significativa de mejora en el desempeño del algoritmo propuesto, para los problemas de optimización''.

    \item Hipótesis alternativa (H1): ``Existe una diferencia significativa de mejora en el desempeño de la técnica propuesta, en el proceso de resolución de problemas de optimización, comparado con el enfoque original del método que no incorporan la optimización del proceso de búsqueda.''.

\end{itemize}


Bajo esta hipótesis H0, se considera la posibilidad de que la propuesta no muestre una diferencia de mejora sobre el algoritmo originario comparado. Esta además implicaría, que cualquier diferencia en la calidad de soluciones entre el algoritmo propuesto y su versión de origen, no se debe a el algoritmo propuesto siendo superior, sino que podría deberse al azar o a otros factores. Por otro lado, en H1 se estaría asumiendo que la propuesta es superior en términos de calidad de solución comparado con el otro algoritmo de referencia.

El par de hipótesis para todo el diseño experimental tiene la finalidad de averiguar a nivel general si: ``Hay diferencia de mejora en la calidad del algoritmo propuesto para los problemas planteados en los experimentos comparado con todas las versiones originales del algoritmo en cuestión''.




\section{Objetivo General}
\label{sc:OG}
Evaluar el impacto del método bioinspirado mejorado, que incorpora técnicas de XAI y estrategias para la detección de estancamientos en óptimos locales, en la calidad de las soluciones obtenidas en problemas de optimización continua y discreta, comparado con la versión original del método, con el fin de determinar si la mejora en la optimización del proceso de búsqueda resulta en un desempeño significativamente superior.



\section{Objetivos Específicos}
\label{sc:OE}

\begin{itemize}
    \item Investigar métodos de explicabilidad de modelos XAI que permitan seleccionar las variables de decisión más relevantes del análisis bayesiano durante el proceso de optimización.
    \item Desarrollar un procedimiento para recopilar y calcular métricas de los datos generados por los algoritmos metaheurísticos durante el proceso de optimización.
    \item Validar la técnica propuesta utilizando problemas de optimización continua CEC LSGO 2023 y CEC 2017, junto a un problema discreta Profitability Tour Problem (PTP).
    \item Evaluar la eficacia de la técnica desarrollada en términos de la calidad de las soluciones obtenidas durante el proceso de optimización, comparándolas con enfoques tradicionales que no incorporan la técnica propuesta, utilizando métricas de rendimiento como el fitness, convergencia, calidad de solución, robustez (análisis estadístico).
    \item Reportar los resultados obtenidos de la propuesta, proporcionando una explicación detallada del comportamiento del algoritmo durante el proceso de optimización.
\end{itemize}
